\section{Related Work}

\subsection{Traditional Light Field Depth Estimation Methods}

Early light field depth estimation methods predominantly rely on handcrafted features and geometric or optical physical models to infer scene structure. These approaches can be broadly categorized into Epipolar Plane Image (EPI) structure analysis, defocus and correspondence cue fusion, and multi-view stereo (MVS) matching.

Wanner and Goldluecke [2] formulated depth estimation as a global optimization problem over EPIs. Their method utilizes the structure tensor to extract the dominant orientation of EPI lines, where the slope is inversely proportional to scene depth. A variational framework is subsequently employed to enforce spatial smoothness and resolve ambiguities. While this approach yields accurate depth maps for Lambertian surfaces with distinct textures, the structure tensor becomes highly sensitive to intensity variations. Specifically, specular highlights introduce abrupt gradient changes, leading to erroneous slope estimations and fragmented depth boundaries.

To address textureless regions where EPI slopes are inherently ambiguous, subsequent research integrated defocus cues with correspondence cues. Tao et al. [1] demonstrated that combining depth from defocus and depth from correspondence significantly improves robustness in weakly textured areas. The defocus cue evaluates the variance of refocused images across different focal planes, while the correspondence cue measures photo-consistency across varying viewpoints. Although this fusion mitigates ambiguity in uniform regions, it inherently assumes that the photo-consistency metric remains invariant across viewpoints. When non-Lambertian reflections are present, the viewpoint-dependent intensity changes violate this invariance, causing severe mismatches in both defocus variance and correspondence costs.

Shifting from continuous EPI analysis to discrete multi-view matching, benchmark frameworks and associated MVS techniques [4] adapted traditional stereo matching to light field data. These methods typically employ shifted photo-consistency metrics or adaptive patch matching to handle the sub-pixel disparities characteristic of light fields. By evaluating the similarity of image patches across multiple sub-aperture images, they achieve high precision in well-textured Lambertian scenes. However, the reliance on fixed-size patch matching makes these methods vulnerable to depth discontinuities and non-Lambertian effects. Specularities shift across patches at different viewpoints, substantially degrading the matching scores and resulting in depth artifacts around reflective boundaries.

Despite their mathematical elegance and effectiveness under ideal conditions, these traditional methods share substantial limitations when applied to complex real-world environments. Primarily, they heavily rely on the Lambertian assumption, positing that surface radiance remains constant regardless of the viewing angle. When surfaces exhibit non-Lambertian characteristics, such as specular or glossy reflections, the continuous EPI slopes fracture or bend, and the photo-consistency metrics fail, rendering traditional geometric features and matching mechanisms ineffective. Second, these approaches lack a physical mechanism to decouple depth from material reflectance. They cannot explicitly perceive or separate different reflection types, leading to highly compromised robustness in non-Lambertian regions. This inherent inability to distinguish between geometric parallax and material-induced intensity variations provides the primary motivation for our three-layer angular signal decomposition physical model. By explicitly isolating ambient lighting, parallax, and material reflectance, the proposed model restores reliable depth estimation in scenarios where traditional geometric assumptions collapse.

\subsection{Deep Learning-based General Light Field Depth Estimation Methods}

Following the limitations of handcrafted features, deep learning techniques have been extensively adopted to extract implicit representations from light field data. Early convolutional neural network (CNN) architectures primarily focus on Epipolar Plane Image (EPI) structure analysis. By treating EPIs as 2D spatial-angular images, these methods employ encoder-decoder structures, such as U-Net [6], to capture the linear slopes corresponding to scene depth. Residual connections [7] are frequently integrated to mitigate gradient vanishing and enhance the extraction of fine-grained EPI textures. These CNN-based approaches demonstrate significant improvements in depth accuracy for Lambertian surfaces by learning robust mapping functions from EPI gradients to disparity values, outperforming traditional structure tensor methods in weakly textured regions.

While EPI-based CNNs excel at local slope extraction, they often struggle with occlusions and view-dependent variations. To address these challenges, subsequent research shifted towards multi-view feature fusion guided by attention mechanisms. Inspired by the success of self-attention in sequence modeling [8], light field networks incorporate channel and spatial attention modules to dynamically re-weight features from different sub-aperture images. This adaptive fusion suppresses inconsistent features caused by occlusions while emphasizing reliable correspondences. Combined with multi-scale feature aggregation, such as Feature Pyramid Networks [9], these attention-driven architectures achieve enhanced robustness in complex scenes by explicitly modeling the reliability of angular observations before depth regression.

More recently, the inherent local receptive field of CNNs has prompted the adoption of Vision Transformers to capture long-range dependencies across the entire light field. Hierarchical architectures, such as the Swin Transformer [10], partition the light field into shifted windows to compute self-attention efficiently, enabling the model to establish global geometric consistency across widely separated viewpoints. By treating angular patches as sequential tokens [11], these Transformer-based methods effectively resolve depth ambiguities in large textureless regions and maintain structural coherence at object boundaries. The global context modeling capability substantially improves depth estimation quality compared to purely local convolutional operations.

Despite the strong feature learning capabilities of these deep architectures, they share critical limitations when applied to real-world mixed-material environments. First, these methods predominantly employ a single, unified network architecture to process all pixels, implicitly relying on a universal Lambertian reflectance assumption. When confronted with mixed scenes containing both Lambertian and Non-Lambertian surfaces, this single EPI assumption generates severe conflicts. Specular reflections and complex textures disrupt the continuous EPI slopes, leading to blurred depth estimations and structural distortions in Non-Lambertian regions. This architectural rigidity highlights the necessity of the dual-mask routing mechanism to decouple and process distinct material domains adaptively. Second, the implicit feature learning in these networks lacks physical interpretability. When trained on highly imbalanced datasets, where Non-Lambertian scene data is extremely scarce, the models tend to overfit to the majority Lambertian domain, resulting in restricted generalization capabilities. This vulnerability to data distribution shifts underscores the core value of the domain-balanced sampling strategy to maximize model robustness under severe data constraints.

\subsection{Deep Learning-based Non-Lambertian and Material-Aware Light Field Methods}

To address the violation of the Lambertian assumption, recent deep learning approaches have introduced dual-branch architectures designed to separate specular and diffuse reflections. These methods typically employ frequency-domain analysis, such as 2D Discrete Fourier Transform (2D-DFT), to isolate high-frequency specular components from low-frequency diffuse signals. Networks incorporating hierarchical feature extraction, like the Swin Transformer [10], have been adapted to process sub-aperture images in parallel branches. One branch focuses on extracting robust diffuse features for depth regression, while the other models the view-dependent specular shifts. By fusing these features through cross-attention modules [8], these architectures achieve notable depth accuracy on moderately reflective surfaces. However, the reliance on 2D-DFT for material separation becomes problematic when the angular resolution is low (e.g., $9 \times 9$ discrete sampling), as frequency aliases severely degrade the isolation of specular highlights.

Building upon dual-branch designs, subsequent research has explored material-aware routing mechanisms to dynamically process different surface types. Instead of explicit physical separation, these methods utilize data-driven mask generation to route pixels into specialized processing streams. Architectures leveraging masked representation learning, such as Masked Autoencoders [12], or segmentation-oriented networks like U-Net [6], are employed to predict implicit material masks. These masks guide the network to apply distinct disparity estimation strategies for Lambertian, specular, and translucent regions. While this implicit routing improves performance on mixed-material scenes, the mask generation relies entirely on supervised learning from limited annotations. The absence of explicit geometric priors makes the mask prediction highly susceptible to overfitting, particularly when training data for non-Lambertian scenes is extremely scarce (e.g., only four scenes available in standard benchmarks [4]).

To mitigate the limitations of purely data-driven routing, recent efforts have integrated physical priors into the material perception process. Prior studies demonstrate that decomposing light field angular signals into DC, linear, and residual components physically isolates depth and material properties [18]. Building on this physical model, geometric angular features, such as local gradients and angular correlations, provide a robust alternative to frequency-domain methods for characterizing surface reflectance [17]. A random forest classifier leveraging these geometric features achieves highly accurate pixel-level discrimination between Lambertian and non-Lambertian regions [19]. Integrating this geometry-based material perception mechanism as a prior mask significantly enhances the reliability of subsequent depth estimation [20]. A three-layer angular signal decomposition model effectively decouples ambient lighting, parallax, and material reflectance, resolving the essential breakdown of Epipolar Plane Image assumptions on non-Lambertian surfaces [21]. Extracting these geometric angular features instead of relying on 2D-DFT effectively bypasses resolution constraints in low-angular-resolution light fields [22].

Despite these advancements, existing material-aware light field methods exhibit critical limitations that hinder their generalization in real-world applications. First, material perception heavily depends on 2D-DFT or similar frequency-domain analyses, which fail under low angular resolutions (e.g., $9 \times 9$ discrete sampling) due to severe spectral aliasing. Current frameworks also lack a strict physical decoupling of angular signals into DC, linear, and residual components, preventing an effective resolution of EPI slope ambiguities. This deficiency directly motivates our adoption of geometric angular features over frequency-domain representations and the design of a three-layer physical model. Second, existing mask generation mechanisms are predominantly implicit and data-driven, lacking efficient explicit priors based on geometric characteristics. Concurrently, current frameworks fail to address the severe overfitting caused by the extreme scarcity of non-Lambertian training data (e.g., only four scenes). The absence of targeted training strategies, such as domain-balanced sampling, restricts their generalization capability. These unresolved issues substantiate the necessity of our unified dual-mask routing framework and the proposed domain-balanced sampling strategy, which collectively maximize depth estimation accuracy across diverse and data-limited scenarios.

Although existing dual-branch methods attempt to separate specular and diffuse reflections, their reliance on frequency-domain analysis for material perception is fundamentally hindered by discrete angular sampling limitations, often leading to EPI slope blurring and poor generalization under scarce non-Lambertian training data. Inspired by these observations, we propose a unified dual-mask physical approach that addresses the aforementioned limitations by introducing a geometry-driven material perception mechanism via three-layer angular signal decomposition for robust routing, and employing a dual-mask framework with domain-balanced sampling to mitigate EPI structural degradation and maximize cross-domain generalization. Building upon these insights, we now formally detail the complete pipeline and architectural design of our proposed framework.